\section{Formal Model}
\label{sec:model}

This section defines the release computation and the publish protocol;
Section~\ref{sec:guarantees} proves guarantees over these definitions.
The model separates stage availability from cross-run admission and discharge.


\subsection{Stage Availability}
For a selected package $p$, let $B_p$ and $P_p$ denote successful completion
of its build and publish stage, respectively. A version task, when present,
is a prefix of $B_p$. For a provider $q$ and consumer $p$, the normal
successful-stage dependencies include $B_q\to B_p$, $B_p\to P_p$, and
$P_q\to P_p$. If $q$ requires published-artifact availability, add
$P_q\to B_p$. Here stage arrows mean prerequisite-to-successor; the package
edge convention below is consumer-to-provider. The construction is over the
full package graph, with stages outside the selected work treated as already
satisfied only when their required artifacts are available.

\begin{proposition}[stage-order safety]
\label{prop:stages}
For a finite acyclic package graph, the constructed stage graph is acyclic.
Any schedule that starts a stage only after all its prerequisites succeed
respects both local-build and published-artifact availability.
\end{proposition}
\begin{proof}
Take a provider-first topological rank of packages. Every cross-package stage
edge increases that rank. Within a package, the only direction is build to
publish. A directed cycle would therefore require a rank decrease or a
publish-to-build edge within one package, neither of which exists. Each
availability requirement is an explicit prerequisite, so a legal schedule
cannot cross it early. The claim concerns ordering, not whether a command
actually produces the promised artifact; that is an integration contract.
\end{proof}

A global build phase followed by a global publish phase violates the
$P_{\mathrm{base}}\to B_{\mathrm{app}}$ edge in Fig.~\ref{fig:stages}.
Conversely, no edge requires $P_{\mathrm{assets}}$ to precede
$B_{\mathrm{base}}$, so those stages may overlap. This elementary construction
uses standard topological scheduling~\cite{kahn,mokhov}; the recovery results
below address the separate problem of retaining the correct selected work.

\subsection{Workspace, State, Windows}

\begin{definition}[history]
The history is a finite DAG of commits with a distinguished
$\mathsf{HEAD} = h$; $\mathit{reach}(c)$ is the set of commits reachable
from $c$ by parent edges, including $c$; write $c' \preceq c$ for
$c' \in \mathit{reach}(c)$. $h$ is fixed throughout. Normalization supplies a deterministic total order
of the commits consistent with ancestry; ``newest'' refers to this order,
with unit position breaking ties within a commit. Ancestry alone is only a
partial order, so this ordering policy is a fixed planning input. This reading of a
repository is the standard one: the commit DAG with reachability is the
working model of git history in repository mining~\cite{birdgit}, and
the same structure underlies formal derivations of git's three-way
merge from least common ancestors~\cite{mrdt}. Unlike the
patch-theoretic accounts of version
control~\cite{swierstraloh,hpt,mimram}, this model tracks no file
contents, only history, tags and normalized intent. Manifests and path-derived
scopes are supplied through the normalized instance rather than modeled as files.
\end{definition}

\begin{definition}[workspace]
$\mathcal{P}$ is the finite set of packages at $h$; $E \subseteq
\mathcal{P} \times \mathcal{P}$ contains $(a,b)$ when $a$ declares a
dependency on $b$ of a configured kind; $\mathit{dep}(p) = \{d : (d,p)
\in E\}$. The graph is acyclic; cyclic workspaces are outside this protocol
and are rejected before scheduling. For $S \subseteq
\mathcal{P}$, $\delta(S,d)$ is the least number of edges on a chain from
a member of $S$ to $d$ against the arrows, $\infty$ if none.
\end{definition}

\begin{definition}[bumps, versions, channels]
$\mathbb{B} = \{\mathsf{none} < \mathsf{patch} < \mathsf{minor} <
\mathsf{major}\}$ is a chain, hence a join-semilattice under
$\sqcup = \max$. Versions are SemVer~2.0.0~\cite{semver} with its
precedence; $\mathit{ch}(v)$ is $\mathsf{stable}$ for a non-prerelease
$v$ and otherwise $v$'s first prerelease identifier; $\mathit{core}(v)$
and $\mathit{cnt}(v)$ are the numeric triple and the prerelease counter;
$\mathit{apply}(v,b)$ applies a bump.
\end{definition}

\begin{definition}[state]
\label{def:state}
A state $\sigma$ is a finite set of release tags $(p, v, c)$ with
$c \preceq h$; states are ordered by inclusion. $\mathit{base}_\sigma(p)$
is the highest-precedence version among $p$'s tags,
$\mathit{sbase}_\sigma(p)$ the highest-precedence \emph{non-prerelease}
one, $\mathit{basec}_\sigma(p)$ and $\mathit{scom}_\sigma(p)$ their
commits; $\mathit{ch}(\mathit{base}_\sigma(p)) = \mathsf{stable}$ by
convention when $p$ has no tag. Equal-precedence tags with conflicting release identities are rejected.
When no stable tag exists, $\mathit{sbase}$ is the virtual 0.0.0 and
$\mathit{reach}(\mathit{scom})=\emptyset$. An untagged package additionally
has $\mathit{base}=0.0.0$, $\mathit{reach}(\mathit{basec})=\emptyset$, and a configured
initial version; applying its first nonzero bump returns that initial version.
Within a valid run sequence, tags are never moved or deleted, so states only
grow. Rewriting a completion record is outside the guarantees.
Tag write access is therefore release authority; a manifest's version field
is not the release ledger. Manifests, configuration, and derived scopes remain
trusted planning inputs.
\end{definition}

\begin{definition}[train and fresh-work windows]
\label{def:window}
\begin{equation}
  W_\sigma(P) \;=\; \mathit{reach}(h) \setminus
                    \mathit{reach}(\mathit{scom}_\sigma(P)),
  \label{eq:window}
\end{equation}
and
\begin{equation}
  W^\circ_\sigma(P) = W_\sigma(P) \setminus
                       \mathit{reach}(\mathit{basec}_\sigma(P)),
  \label{eq:fresh}
\end{equation}
with an empty subtracted reach when $P$ has no tag.  $W_\sigma$ is the
\emph{train window}; $W^\circ_\sigma$ is the \emph{fresh-work window}.
When $P$ has no stable baseline, $W_\sigma(P) = \mathit{reach}(h)$:
an unreleased package's first release accounts for its whole
history. An existing project must therefore configure its initial baseline
and scope mapping deliberately. The window is measured from the last \emph{stable} tag
so that a prerelease train can recompute its target over the whole train
(Eq.~\ref{eq:train}). Direct admission, propagation admission and ledger
discharge use $W^\circ$ instead. Thus a prerelease tag discharges delivered
work without erasing the earlier records needed to recompute the train target.
\end{definition}

\subsection{Units, Suppression, Admission}

\begin{definition}[units]
$\mathcal{U}$ is the set of units carried by commits in
$\bigcup_P W_\sigma(P)$, parsed and with scopes resolved and corrections
applied. For $u \in \mathcal{U}$: $c_u$ is its commit,
$\mathit{scope}(u) \subseteq \mathcal{P}$ its resolved scope,
$\mathit{bump}(u) \in \mathbb{B}$ the bump its type and \code{!} map to;
its bump axis is $\mathit{pb}(u) \in \mathbb{B}$ (the stated value, or
$\mathit{bump}(u)$ under \code{inherit}; the default once propagation is
asked for is $\mathsf{patch}$ because a dependency update need not
change the consumer's own API), $\mathit{pd}(u) \in \mathbb{N} \cup
\{\infty\}$ (default $0$), $\mathit{ps}(u)$ a scope filter; its channel
axis is a channel value $\mathit{pc}(u)$, a depth $\mathit{pcd}(u)$
(default $0$) and a filter $\mathit{pcs}(u)$.
\end{definition}

\begin{definition}[suppression and sources]
\label{def:src}
$\mathit{canc}(c,p)$ holds when $c$ is an ancestor-or-self of a
\code{cancel} commit whose scope contains $p$; $\mathit{held}(p)$ when
$p$'s effective \code{Release-As} is \code{none}. Both are functions of
$h$ and the configuration alone. With $\mathit{disch}_\sigma(p,c)
\equiv c \notin W^\circ_\sigma(p)$,
\begin{equation}
\begin{aligned}
  \mathit{src}_\sigma(u) = \{ p \in \mathit{scope}(u) : {}&
    \mathit{disch}_\sigma(p,c_u) \\
    &\lor\lnot(\mathit{canc}(c_u,p) \lor \mathit{held}(p)) \}.
\end{aligned}
  \label{eq:src}
\end{equation}
Suppression applies only to \emph{undischarged} work: once $p$ has
published the version carrying $u$, the artefact its consumers are owed
is public, and neither a later cancel nor a later hold may retroactively
strand them.
\end{definition}

\begin{definition}[reach]
\label{def:reach}
$R(S,\ell) = \{ d \in \mathcal{P} \setminus S : \delta(S,d) \leq \ell
\}$. Excluding all of $S$ means a unit never propagates to a package it
bumps directly; depth is shortest-path and always measured from the
originating source set, and a propagated effect never cascades its own
propagation, so the depth written in the message is the whole blast
radius in every run.
\end{definition}

\begin{definition}[channels]
\label{def:chan}
$\mathit{chan}_\sigma(P)$, resolved for \emph{every} package: the value
proposed by the newest eligible direct \code{Channel} directive naming $P$ with
$c_u \in W^\circ_\sigma(P)$ and $\lnot\mathit{canc}(c_u,P)$; else the value proposed for $P$ by the newest eligible
channel-axis unit $u$ with $\mathit{pcd}(u) > 0$, $P \in
R(\mathit{src}_\sigma(u), \mathit{pcd}(u)) \cap \mathit{pcs}(u)$,
$c_u \in W^\circ_\sigma(P)$ and $\lnot\mathit{canc}(c_u,P)$; else
$\mathit{ch}(\mathit{base}_\sigma(P))$. A proposal is computed against
$\mathit{cur} = \mathit{ch}(\mathit{base}_\sigma(P))$ (the tag, never
a value computed this run) and is empty when it proposes
$\mathit{cur}$ itself; a transition proposes only where $\mathit{cur}$
matches its \code{<from>}; a \emph{propagated} non-transition
$\mathsf{stable}$ never graduates a prerelease package (graduation must
be deliberate); \code{inherit} resolves to the unit's own
\code{Channel} value if any, else to
$\mathit{ch}(\mathit{base}_\sigma(q))$ for its sources $q$, bytewise
least name deciding a disagreement.
An empty proposal is ineligible; resolution chooses the newest remaining
proposal, with direct proposals taking precedence over propagated ones.
\end{definition}

\begin{definition}[admission and effective bump]
\label{def:adm}
Let $\mathit{live}(u)$ mean $\mathit{bump}(u)\neq\mathsf{none}$,
$\mathit{pb}(u)\neq\mathsf{none}$, and $\mathit{pd}(u)>0$. With
$\mathit{srcCh}_\sigma(u)$ equal to
$\{\mathit{chan}_\sigma(p) : p \in \mathit{src}_\sigma(u)\}$ and
\begin{equation}
  \mathit{rsv}_\sigma(u,D) \equiv \mathsf{stable} \in
  \mathit{srcCh}_\sigma(u) \lor \mathit{chan}_\sigma(D) \in
  \mathit{srcCh}_\sigma(u)
  \label{eq:resolvable}
\end{equation}
(the modeled policy admits a prerelease source only on a matching
consumer channel; this is an explicit release policy, not a universal
claim about package installers),
\begin{equation}
\label{eq:admission}
\begin{aligned}
  \mathit{adm}_\sigma(u,D) \;\equiv\; & \\[-0pt]
     & D \in R(\mathit{src}_\sigma(u), \mathit{pd}(u))
     &\!\!& \text{(reached)} \\
  \wedge\;\; & D \in \mathit{ps}(u)
     &\!\!& \text{(in scope)} \\
  \wedge\;\; & c_u \in W^\circ_\sigma(D)
     &\!\!& \text{(fresh window)} \\
  \wedge\;\; & \neg\,\mathit{canc}(c_u, D)
     &\!\!& \text{(not discarded)} \\
  \wedge\;\; & \mathit{live}(u)
     &\!\!& \text{(bump-live)} \\
  \wedge\;\; & \mathit{rsv}_\sigma(u, D)
     &\!\!& \text{(resolvable),}
\end{aligned}
\end{equation}
\begin{equation}
\label{eq:effective}
  \mathit{eff}_\sigma(P) = \mathit{dir}_\sigma(P) \sqcup
  \bigsqcup \{\, \mathit{pb}(u) : \mathit{adm}_\sigma(u,P) \,\},
\end{equation}
where $\mathit{dir}_\sigma(P) = \bigsqcup\{\mathit{bump}(u) : P \in
\mathit{src}_\sigma(u) \cap \mathit{scope}(u),\ c_u \in W^\circ_\sigma(P)\}$.
Because $\sqcup$ is a join, a package with its own \code{feat} pending
is unaffected by an incoming propagated \code{patch}, joining an admitted bump cannot lower the effective bump, and gathering order cannot matter.
\end{definition}

The consumer-side staleness audit is the upward dual of
Eq.~\ref{eq:admission}, not merely reverse reachability. For candidate $D$
it scans units in $W^\circ_\sigma(D)$, excludes $D$ when it is itself a
source of that unit, and applies the same shortest depth, scope,
cancellation, $\mathit{live}$, and resolved-channel predicate $\mathit{rsv}$. These clauses
make the upward answer exactly equal to the downward planner's obligations,
including prerelease workspaces.

\subsection{Plan, Versions, Runs}

\begin{definition}[plan and versions]
\label{def:plan}
Write $\mathit{pin}_\sigma(P)=v$ when the effective surviving exact
\code{Release-As: v} unit for $P$ has $c_u\in W^\circ_\sigma(P)$; otherwise
$\mathit{pin}_\sigma(P)$ is undefined. Thus a successful tag discharges a
pin just like other fresh work, and the directive is fixed during a retry.
\begin{equation}
\begin{aligned}
  \Pi(\sigma) = \{\, P :\;
    &\lnot\mathit{held}(P) \wedge (\mathit{eff}_\sigma(P) \neq
     \mathsf{none} \\
    &\ \lor\ \mathit{chan}_\sigma(P) \neq
     \mathit{ch}(\mathit{base}_\sigma(P))
     \\
    &\ \lor\ \mathit{pin}_\sigma(P)\downarrow) \,\},
\end{aligned}
\end{equation}
each member carrying $\nu_\sigma(P)$. Let $S=\mathit{sbase}_\sigma(P)$,
$B=\mathit{base}_\sigma(P)$. Define the train aggregate
$\mathit{train}_\sigma(P)$ by Eq.~\ref{eq:effective}, replacing its direct
and propagated fresh-window tests $c_u\in W^\circ_\sigma(P)$ with
$c_u\in W_\sigma(P)$. All other source, scope, suppression, depth, and
resolved-channel predicates stay the same. It computes the target bump over
the train, while $\mathit{eff}$ determines fresh admission. Write
$e=\mathit{train}_\sigma(P)$.
If $\mathit{pin}_\sigma(P)\downarrow$, it supplies that exact version. A stable package that
remains stable uses $\mathit{apply}(S,\mathit{eff}_\sigma(P))$; a prerelease
package graduating to stable uses $\mathit{apply}(S,e)$ as required by
\lang{}~2.0 \S11.5. Otherwise, with $\mathit{ch}$ the target channel,
\begin{equation}
\label{eq:train}
\begin{aligned}
  t &= \mathit{apply}(S, e), \\
  \nu_\sigma(P) &=
  \begin{cases}
    t\text{-}\mathit{ch}.(\mathit{cnt}(B){+}1)
      & \begin{array}{l}\text{if } \mathit{core}(B) = t\text{ and}\\
        \mathit{ch}(B) = \mathit{ch},\end{array}\\
    t\text{-}\mathit{ch}.0 & \text{otherwise,}
  \end{cases}
\end{aligned}
\end{equation}
For a channel-only release whose candidate does not exceed $B$, recompute
once with $e=\mathsf{patch}$ (the channel-entry patch); an exact pin is
never adjusted this way. This permits entry from a clean stable baseline,
for example 1.2.0 to 1.2.1-beta.0, without treating an actual bump regression
as a channel switch. Every computed or pinned version must exceed
$\mathit{base}_\sigma(P)$ in precedence, or the run refuses. Recomputing
$t$ from the \emph{stable} baseline every run is what moves a whole
train to \code{2.0.0-beta.0} when a breaking change lands mid-train
(Table~\ref{tab:train}).

\end{definition}

\begin{table}[t]
\caption{Eq.~\ref{eq:train} tracing one train.}
\label{tab:train}
\small
\begin{tabular}{@{}lll@{}}
\toprule
Baseline & Pending bump, channel & Next \\
\midrule
\code{1.2.3}        & minor, \code{beta}          & \code{1.3.0-beta.0} \\
\code{1.3.0-beta.0} & no fresh commit             & empty plan \\
\code{1.3.0-beta.0} & a fresh minor commit        & \code{1.3.0-beta.1} \\
\code{1.3.0-beta.1} & minor $+$ a new fix         & \code{1.3.0-beta.2} \\
\code{1.3.0-beta.2} & a breaking change lands     & \code{2.0.0-beta.0} \\
\code{1.3.0-beta.2} & none, switch to \code{rc}   & \code{1.3.0-rc.0} \\
\code{1.3.0-rc.1}   & none, \code{stable}         & \code{1.3.0} \\
\bottomrule
\end{tabular}
\end{table}

\begin{definition}[run]
\label{def:run}
A run from $\sigma$ is parameterised by an adversarially chosen set
$F \subseteq \Pi(\sigma)$ of packages failing before recorded completion, which induces the
blocked set
$B = \{p \in \Pi(\sigma) \setminus F : p$ transitively depends on some
$q \in F\}$, computed over the \emph{full} workspace graph (a package
with no bump this run is still a path from a dependent to a failed
dependency). Publication proceeds in a topological order of the full
graph, ties broken bytewise by name, filtered to the plan \emph{last}:
inducing the graph on the plan first deletes ordering paths that run
through unplanned packages. In the abstract transition, each package is tagged after its
own publish succeeds, never before (a tag for an unpublished version
makes it unrecoverable, since Eq.~\ref{eq:window} treats it as
released), and \emph{whatever its publish target},
including targets that upload nothing
(Proposition~\ref{prop:tagall}). The run yields
$\mathit{run}(\sigma,F) = \sigma \cup \{(p, \nu_\sigma(p), h) : p \in
\Pi(\sigma) \setminus (F \cup B)\}$. The transition abstracts the publish-to-record interval; it does not
prove that interval atomic in the implementation. A \emph{run sequence} is
$\sigma_1 \leq \sigma_2 \leq \cdots$ with $\sigma_{k+1} =
\mathit{run}(\sigma_k, F_k)$.
\end{definition}

The propagation bound excludes parsing, history traversal, scope resolution,
and output. Those costs depend additionally on reachable commits, parent
edges, and retained provenance. Cached traversal may share reachability, but
must preserve each unit's source set, depth, filters, channel predicates, and
identity. Merging units by directive spelling alone can change admission.

Order matters for a reason tags cannot express: a dependent can be published
while its promised provider version is unavailable. If both completion tags
later point at the same commit, ancestry alone does not reveal that temporal
inversion. The protocol therefore prevents it through stage prerequisites. When a
publish is rejected because the version already exists (a run died
between publish and tag), \lang{}~2.0 requires content-identity verification
before adoption and otherwise aborts: tagging a foreign artefact would enrol
it into the release lineage. This is a protocol obligation, not a claim that
\tool{} can inspect arbitrary user-defined publish scripts; such integrations
must supply an idempotency or identity check.

\begin{theorem}[determinism]
\label{thm:det}
$\Pi(\sigma)$, $\nu_\sigma$, the canonical publish order, and the
version-guard refusal outcome are functions of the normalized instance $I$
defined below, including its fixed history-ordering policy.
\end{theorem}
\begin{proof}
Every definition above is a closed expression in those inputs; the three
choice points (competing directives, competing propagated channels,
publish-order ties) are each resolved by a stated total order: newest
commit then last unit within it, and bytewise least name.
\end{proof}

\subsection{Instantiation: What the Theorems Say About a Repository}
\label{sec:model:inst}

A normalized instance is
\[
 I=(\text{history},h,\mathcal P,E,\mathcal U,\sigma,\mathit{config}),
\]
where configuration includes the fixed history-ordering policy. Repository
files and messages supply these objects through normalization. The model
abstracts artifact contents, arbitrary hook effects, and other configuration
features. Its guarantees quantify over this fragment under acyclicity,
version validation, and the run-sequence retry invariant. The implementation
is not assumed to check every component of that invariant. Proving that a
binary refines the model remains separate work; Section~\ref{sec:eval}
provides tests of selected behavior, not a refinement proof.
